% fuzz -t docs/peripheral_registry_lifecycle.tex
% PROBCLI=/opt/prob/probcli bash scripts/verify_peripheral_registry_model.sh
\documentclass[a4paper,10pt,fleqn]{article}
\usepackage[margin=1in]{geometry}
\usepackage{fuzz}
\begin{document}
\title{Hub Peripheral Registry Lifecycle}
\author{A bounded model of static endpoint and worker composition}
\date{}
\maketitle

\section{Boundary and refinement}
The model represents two peripheral kinds, not an audio protocol. Each kind
abstracts its endpoint's opaque key; exact key identity and port typing are
checked separately by the mirrored key and protocol tests. Registration and
binding are separate sequences, so workers can bind in a different order.
The bound of two entries includes a duplicate kind in the uniqueness mutant.
Safety properties are reachability goals, not state predicates that could
silently disable incorrect transitions.

\texttt{Register} refines \texttt{PeripheralRegistry.register};
\texttt{Bind} refines \texttt{bind\_workers}. \texttt{Start} and \texttt{Stop}
validate complete binding before sealing, then record their dispatch plan.
An implementation worker can raise and truncate that plan to a prefix; there
is no retry or rollback. Worker completion, thread scheduling and Display
transport are outside this composition model. Repeated start/stop, including
stop before start and an empty registry, remain allowed as in the implementation.
Rejected operations have no transition: their error types are unit-tested.

\section{Types and fidelity switches}
\begin{zed}
  ZBOOL ::= ztrue | zfalse \\
  RGKIND ::= rkdisplay | rkaudio \\
  RGPHASE ::= rgopen | rgsealed \\
  RGACTION ::= rgnone | rgregister | rgbind | rgstart | rgstop
\end{zed}
\begin{axdef}
  UNIQUE, REGISTERED, FROZEN, COMPLETE, STARTORDER, STOPORDER : ZBOOL
\where
  UNIQUE = ztrue \\
  REGISTERED = ztrue \\
  FROZEN = ztrue \\
  COMPLETE = ztrue \\
  STARTORDER = ztrue \\
  STOPORDER = ztrue
\end{axdef}

\section{State and operations}
\texttt{previous} and \texttt{action} are observational ghosts. The no-mutation
goal checks the pre-state of register/bind, not merely a sealed post-state.
\begin{schema}{Registry}
  registration, bindings, dispatch : \seq RGKIND \\
  phase, previous : RGPHASE \\
  action : RGACTION
\where
  \# registration \leq 2 \\
  \# bindings \leq 2 \\
  \# dispatch \leq 2
\end{schema}
\begin{schema}{Init}
  Registry
\where
  registration = \langle \rangle \land bindings = \langle \rangle \\
  dispatch = \langle \rangle \\
  phase = rgopen \land previous = rgopen \land action = rgnone
\end{schema}
\begin{schema}{Register}
  \Delta Registry \\
  kind? : RGKIND
\where
  FROZEN = zfalse \lor phase = rgopen \\
  \# registration < 2 \\
  UNIQUE = zfalse \lor kind? \notin \ran registration \\
  registration' = registration \cat \langle kind? \rangle \\
  bindings' = bindings \land dispatch' = dispatch \\
  phase' = phase \land previous' = phase \land action' = rgregister
\end{schema}
\begin{schema}{Bind}
  \Delta Registry \\
  kind? : RGKIND
\where
  FROZEN = zfalse \lor phase = rgopen \\
  \# bindings < 2 \land kind? \notin \ran bindings \\
  REGISTERED = zfalse \lor kind? \in \ran registration \\
  bindings' = bindings \cat \langle kind? \rangle \\
  registration' = registration \land dispatch' = dispatch \\
  phase' = phase \land previous' = phase \land action' = rgbind
\end{schema}
\begin{schema}{Start}
  \Delta Registry
\where
  COMPLETE = zfalse \lor \ran bindings = \ran registration \\
  phase' = rgsealed \land previous' = phase \land action' = rgstart \\
  registration' = registration \land bindings' = bindings \\
  (STARTORDER = ztrue \land dispatch' = registration) \lor
  \quad~ (STARTORDER = zfalse \land dispatch' = bindings)
\end{schema}
\begin{schema}{Stop}
  \Delta Registry
\where
  COMPLETE = zfalse \lor \ran bindings = \ran registration \\
  phase' = rgsealed \land previous' = phase \land action' = rgstop \\
  registration' = registration \land bindings' = bindings \\
  (STOPORDER = ztrue \land dispatch' =
  \quad~ \{ i : 1 \upto \# registration @
    i \mapsto registration(\# registration + 1 - i) \}) \lor
  \quad~ (STOPORDER = zfalse \land dispatch' = registration)
\end{schema}

\section{Verification and partitions}
The gate exhaustively excludes duplicate registered kinds, bindings without
endpoints, register/bind from a sealed pre-state, incomplete lifecycle dispatch,
start in anything but registration order, and stop in anything but reverse
registration order. It separately checks deadlock freedom: an open registry
can bind a missing worker; a completely bound registry can start or stop.
There is no synthetic escape operation.

Each switch has a defect witness: duplicate registration, unregistered binding,
post-seal registration, incomplete start, incomplete stop, start in binding
order, and stop in forward order. Disabling the bind seal guard alone cannot
enable a successful bind in a reachable sealed state: complete binding and
duplicate-binding rejection also prevent it. The explicit guard is retained,
and its wrong-state rejection is covered directly by unit tests.

The partitions in \texttt{tests/hub/test\_peripheral\_registry.py} cover missing
and exact-key lookup; duplicate kinds with the same or different capability;
unknown and duplicate worker binding; incomplete start/stop before any worker
runs; frozen registration/binding after either lifecycle operation; reverse
binding order; and propagated worker exceptions.
\texttt{tests/hub/test\_display\_workers.py} covers the real Display facade's
four-event trace. Empty and repeated lifecycle dispatch are covered by the
model; the empty case is also unit-tested. This is a bounded composition proof,
not a claim about arbitrary transport behavior or worker thread completion.
\end{document}
